\section{A sparse perturbation of the outer coloring}
\label{perturb:section}

Fix the outer coloring $c_*$ supplied by \Cref{outer:coloring}, and put
$c_0(n)=c_*(L(n))$.  All randomness used to construct $c_*$ has now been
resolved.  We shall introduce new independent bits indexed by geometric
keys.  A key will occur only a bounded number of times on any sampled
progression.  This allows a small flip probability to preserve the
progressions containing many of each outer color.  For progressions with
affine lifts, a separate count of the possible keys makes the same small
probability sufficient.

Retain the parameters and cyclic progressions from
\Cref{setup:parameters,setup:cyclic}.  For $n\in G$, define
\begin{equation}\label{perturb:keys}
 f(n)=\norm{q\tilde y(n)}_2^2,
 \qquad
 \kappa(n)=\left((V(y_i(n)))_{i=1}^D,\,\floor{f(n)}\right),
 \qquad
 \cK=\{\kappa(n):n\in G\}.
\end{equation}
Thus a key records a specific second-system box and a specific integer
squared-norm band.  Since $G$ is finite, so is $\cK$.
The next lemma uses a bounded-width version of the equal-norm geometry
in Behrend's construction~\cite{behrend1946}: a thin norm band has
short intersections with a line on either side of its nearest point
to the origin.

\begin{lemma}[Bounded multiplicity of a key]\label{perturb:multiplicity}
For every progression $n_j=n_0+jd$ in $G$ with $d\ne0$ and
$0\le j<k$, every key occurs at most four times among the $\kappa(n_j)$.
\end{lemma}

\begin{proof}
First fix a second-system box, and let $J$ be the indices whose centered
representatives lie in that box.  Every coordinate interval has width at
most $2H$, so \Cref{dichotomy:affinity} gives
\[
 \tilde y(n_j)=A+jB\qquad(j\in J)
\]
for suitable $A,B\in\R^D$.  The indices in $J$ need not be consecutive.
In particular, the points $q\tilde y(n_j)$ lie on an affine line when
$|J|\ge2$; when $|J|\le1$ the claimed multiplicity bound is immediate.

Any two distinct terms of the progression give points separated by
Euclidean distance at least one after this scaling.  To see this, the
terms are distinct by \Cref{setup:cyclic}, and $\lambda$ is a unit modulo
$q^D$.  Hence $\lambda(n_j-n_{j'})$ is a nonzero residue modulo $q^D$.
Choose an integer representative $z$ of this residue, and let
$a\in\{0,\ldots,D-1\}$ be the largest integer for which $q^a\mid z$.
Divisibility by $q^b$, for $b\le D$, is unchanged upon adding a multiple
of $q^D$, so this choice is intrinsic to the residue.  Write $z=q^a u$;
then $q\nmid u$.  In coordinate $a+1$ the circle displacement is
\[
 y_{a+1}(n_j)-y_{a+1}(n_{j'})=\frac{u}{q}\pmod1.
\]
The residue of $u$ modulo $q$ is one of $1,\ldots,q-1$, whose distance
from $0$ on the circle after division by $q$ is at least $1/q$.
This argument requires only $q\nmid u$; it does not require $u$ to be a
unit or $q$ to be prime.  The difference between any two real
representatives has absolute value at least their circle distance.
Consequently
\begin{equation}\label{perturb:separation}
 \norm{q\tilde y(n_j)-q\tilde y(n_{j'})}_2\ge1
 \qquad(j\ne j').
\end{equation}

If $|J|\ge2$, these separated points determine a nonconstant affine
line $\ell$.  Let $z_0$ be the point of $\ell$ closest to the origin,
and choose a unit vector $e$ along $\ell$.  Writing its points as
$z_0+te$, with $t\in\R$, gives
\[
 \norm{z_0+te}_2^2=t^2+C',\qquad C'=\norm{z_0}_2^2\ge0,
\]
because $z_0$ is orthogonal to $e$.  For a fixed integer $m\ge0$, the
band of the key is
\[
 m\le t^2+C'<m+1.
\]
If it is nonempty, then $m+1-C'>0$, and on the half-line $t\ge0$ it
is contained in the closed interval with endpoints
\[
 \sqrt{\max\{0,m-C'\}},\qquad \sqrt{m+1-C'}.
\]
The difference of the quantities under these square roots is at most
one.  For $0\le u\le v$, the inequality
$(\sqrt v-\sqrt u)^2\le v-u$ therefore shows that this interval has
length at most one.  Reflection gives the same bound on $t\le0$.
By \Cref{perturb:separation}, each of these intervals contains at most
two sampled points: three points in an interval of length at most one
would have total separation at least two.  Thus the key occurs at most
four times.  If the band reaches the nearest point $z_0$, counting
$t=0$ in both halves only enlarges this upper bound.  Replacing the
half-open band by closed intervals also covers equality at either
boundary, and the case of exactly two points needs no additional
argument.
\end{proof}

\begin{figure}[htbp]
 \centering
 \begin{tikzpicture}[scale=1.05,>=Latex,
    every node/.style={font=\small}]
  \path[fill=blue!9,even odd rule]
    (0,0) circle[radius={sqrt(2)}] (0,0) circle[radius=1];
  \draw[gray!75] (0,0) circle[radius=1];
  \draw[gray!75,dashed] (0,0) circle[radius={sqrt(2)}];
  \draw[gray!75,->] (-2.0,0.6)--(2.12,0.6);
  \node[above left] at (-1.82,0.6) {$\ell$};
  \node[above] at (2.02,0.6) {$t$};
  \draw[blue!70!black,line width=2pt]
    ({-sqrt(1.64)},0.6)--(-0.8,0.6);
  \draw[blue!70!black,line width=2pt]
    (0.8,0.6)--({sqrt(1.64)},0.6);
  \fill[blue!70!black] (-0.8,0.6) circle[radius=1.6pt];
  \fill[blue!70!black] (0.8,0.6) circle[radius=1.6pt];
  \draw[blue!70!black,fill=white,line width=0.7pt]
    ({-sqrt(1.64)},0.6) circle[radius=1.6pt];
  \draw[blue!70!black,fill=white,line width=0.7pt]
    ({sqrt(1.64)},0.6) circle[radius=1.6pt];
  \draw[densely dotted] (0,0)--(0,0.6);
  \draw[gray!75] (0,0.46)--(0.14,0.46)--(0.14,0.6);
  \fill (0,0) circle[radius=1pt];
  \node[below left] at (0,0) {$0$};
  \node[above] at (0,0.6) {$z_0$};
  \node[align=left,anchor=west,text width=6.1cm] at (2.65,0.1)
    {$\norm{z_0+te}_2^2=t^2+C'$\par\medskip
     Each half-line contributes a piece\par
     of length at most $1$.\par\medskip
     Sampled points are at least $1$ apart:\par
     at most $2$ points per piece.};
\end{tikzpicture}

 \caption{A line meeting the band $m\le\norm{z}_2^2<m+1$, illustrated
 for $m=1$.  Solid endpoints belong to the inner boundary; open endpoints
 lie on the excluded outer boundary.  In general each half-line from
 the nearest point $z_0$ contributes length at most one, giving the
 four-point bound in \Cref{perturb:multiplicity}.}
 \label{perturb:norm-band}
\end{figure}

Call a cyclic progression \emph{affine} when its entire sequence of
centered representatives $\tilde y(n_j)$ is affine in $j$ as a sequence
in $\R^D$.  Its \emph{signature} is the literal word
\begin{equation}\label{perturb:signature}
 \sigma=\bigl((c_0(n_j),\kappa(n_j))\bigr)_{0\le j<k}.
\end{equation}
In particular, the entries retain the identities of the boxes and the
integer band indices.  Keys are not renamed when passing from one
progression to another.

\begin{lemma}[Counting affine signatures]\label{perturb:signatures}
Let $T_{\mathrm{aff}}$ be the number of signatures realized by affine
cyclic progressions, and let $T_{\mathrm{glob}}$ be the full-label word
count of \Cref{count:global}.  There is an absolute constant $C$ such
that
\begin{equation}\label{perturb:signature-count}
 T_{\mathrm{aff}}
 \le C T_{\mathrm{glob}}\bigl(k(Dq^2+2)\bigr)^3.
\end{equation}
In particular,
\begin{equation}\label{perturb:entropy}
 \log\max\{1,T_{\mathrm{aff}}\}
 \le O(k^{3/10}\log k)+6\log q+O(\log k)
 =O(k^{9/10}\log k)=o(pk).
\end{equation}
All implied constants are absolute.
\end{lemma}

\begin{proof}
On an affine progression write $\tilde y(n_j)=A+jB$.  Then
\[
 f(n_j)=q^2\norm{A+jB}_2^2
       =\theta_0+\theta_1j+\theta_2j^2,
\]
where
$\theta_0=q^2\norm A_2^2$,
$\theta_1=2q^2\langle A,B\rangle$, and
$\theta_2=q^2\norm B_2^2$.
Moreover, since every coordinate of $\tilde y(n_j)$ belongs to
$[-1/2,1/2)$,
\[
 0\le f(n_j)\le Dq^2/4\qquad(0\le j<k).
\]
Put $B_*=\floor{Dq^2/4}+1$.  For each $j\in\{0,\ldots,k-1\}$ and
$b\in\{0,\ldots,B_*\}$, compare
$\theta_0+\theta_1j+\theta_2j^2$ with $b$.  These comparisons use at
most $k(B_*+1)\le k(Dq^2+2)$ affine hyperplanes in the three variables
$(\theta_0,\theta_1,\theta_2)$.  Their joint signs, including equality,
determine the complete word $(\floor{f(n_j)})_{j<k}$ for all realized
progressions.  \Cref{count:arrangement}, in the fixed dimension three,
therefore gives at most $C\bigl(k(Dq^2+2)\bigr)^3$ such words.  Allowing
all triples of coefficients, without imposing their geometric
relations, can only increase this count.

A full-label word $(L(n_j))_{j<k}$ determines the specific second-system
boxes and, because $c_*$ is fixed, every outer bit $c_0(n_j)$.  Together
with the integer-part word just counted, it determines every entry of
\Cref{perturb:signature}.  Multiplying the two counts proves
\Cref{perturb:signature-count}; no independence between the two words
is asserted or needed.

For the logarithmic estimate, \Cref{count:global} bounds
$\log T_{\mathrm{glob}}$ by $O(k^{3/10}\log k)$.  Also
$Dq^2+2\le3Dq^2$, so the remaining logarithm in
\Cref{perturb:signature-count} is at most
$6\log q+3\log k+3\log D+O(1)$.
\Cref{setup:parameters,setup:q-bounds} give
\[
 \log q\le \frac{ck}{D}\log k+\log(2k^2)
          =O(k^{9/10}\log k),
 \qquad pk=k^{19/20}.
\]
Since $k^{9/10}\log k/k^{19/20}=k^{-1/20}\log k\to0$,
\Cref{perturb:entropy} follows.  The arrangement constant comes from
dimension three, so it is independent of the growing dimension $D$.
\end{proof}

\begin{theorem}[A cyclic two-coloring]\label{perturb:cyclic}
There is an absolute integer $K_0$ such that, for every integer
$k\ge K_0$, the construction in \Cref{sec:setup} gives an integer
\[
 N=q^D\ge k^{ck},\qquad c=10^{-5},
\]
and a coloring $C:\Z/N\Z\to\{0,1\}$ for which no progression
$n_0,n_0+d,\ldots,n_0+(k-1)d$ with $d\ne0$ in $\Z/N\Z$ is
monochromatic.
\end{theorem}

\begin{proof}
For each key $\kappa\in\cK$, take an independent Bernoulli bit
$E_\kappa$ with $\Pp(E_\kappa=1)=p=k^{-1/20}$, and define
\begin{equation}\label{perturb:final-coloring}
 C(n)=c_0(n)+E_{\kappa(n)}\pmod2.
\end{equation}
Only this finite family of bits is random.  In particular, positions
with a common key use the very same bit, even if they have different
outer colors.  We may take $k$ sufficiently large that $p<1/2$.

First consider a progression containing at least $\gamma k$ of each
outer color, where $\gamma=1/100$.  Fix a target color $b\in\{0,1\}$.
Every position with $c_0(n_j)=1-b$ must have $E_{\kappa(n_j)}=1$ for
the whole progression to receive color $b$.  By
\Cref{perturb:multiplicity}, these positions specify at least
$\gamma k/4$ distinct keys.  This set of keys is determined before
the bits are sampled.  Thus independence between distinct keys gives
\[
 \Pp\bigl(C(n_j)=b\text{ for all }j<k\bigr)
 \le p^{\gamma k/4}.
\]
If some key also occurs at a position with outer color $b$, the
prescriptions for that key conflict and this probability is zero,
which obeys the same bound.  There is no choice of a subset of keys to
count: all keys requiring a flip are already specified by the
progression and the target color.

There are fewer than $|G|^2=q^{2D}$ choices of $(n_0,d)$ with $d\ne0$.
A union bound over these choices and the two target colors bounds the
probability $R_k$ that any such progression becomes monochromatic by
\begin{equation}\label{perturb:rich-bound}
 \begin{aligned}
 R_k&\le 2q^{2D}p^{\gamma k/4}\\
 &\le\exp\!\left(
       \left(2c-\frac{\gamma}{80}\right)k\log k
       +2D\log(2k^2)+\log2\right).
 \end{aligned}
\end{equation}
Here the coefficient is explicitly
\begin{equation}\label{perturb:negative-coefficient}
 2c-\frac{\gamma}{80}
 =\frac{1}{50000}-\frac{1}{8000}
 =-\frac{21}{200000}<0.
\end{equation}
As $D=O(k^{1/10})$, the other terms in the exponent are
$O(k^{1/10}\log k)=o(k\log k)$.  Hence $R_k=o(1)$.

Next fix a signature realized by an affine progression, and again fix
the target color $b$.  For each entry $(c_j,\kappa_j)$ of this
signature, monochromaticity prescribes
\[
 E_{\kappa_j}=b-c_j\pmod2.
\]
Inconsistent prescriptions give an empty event.  Otherwise, if there
are $s$ distinct keys and $u$ of them are prescribed the value one,
independence gives its probability exactly
$p^u(1-p)^{s-u}$.  By \Cref{perturb:multiplicity}, $s\ge k/4$, and
since $p\le1-p$ this is at most
\[
 (1-p)^s\le(1-p)^{k/4}.
\]
Two progressions with the same literal signature impose precisely the
same prescriptions on the same random variables.  They therefore
define the same event for this target color and need be counted only
once.  Consequently the probability $A_k$ of a monochromatic affine
progression satisfies
\begin{equation}\label{perturb:affine-bound}
 \begin{aligned}
 A_k&\le2T_{\mathrm{aff}}(1-p)^{k/4}\\
 &\le2\max\{1,T_{\mathrm{aff}}\}\exp(-pk/4)
   =\exp\bigl(-pk/4+o(pk)\bigr)=o(1),
 \end{aligned}
\end{equation}
using \Cref{perturb:signatures} and $1-p\le e^{-p}$.
Neither union bound requires independence between different
progressions or signatures.

By \Cref{dichotomy:main}, every nonzero-step cyclic progression is
covered by the first case or the affine case.  The cases may overlap,
which only enlarges their combined bound.  Thus the probability that
any such progression is monochromatic is at most $R_k+A_k=o(1)$.
For all sufficiently large $k$ this sum is less than one, so a choice
of the bits gives the asserted coloring.  All preceding thresholds and
constants depend only on the fixed numerical parameters, and the
bound in \Cref{setup:q-bounds} is uniform over every prime $P$
allowed by \Cref{setup:prime}. A single absolute $K_0$ therefore works for
every integer $k\ge K_0$.  Finally, $q\ge k^{ck/D}$ gives
$N=q^D\ge k^{ck}$.
\end{proof}
